The statement in the problem uses asymptotic notation ($\gg$) and does not fully specify dependencies when $\delta=\delta(n)$.
The minimal conventional restatement consistent with the text is:

For each $n\ge 1$ and each real $\delta\in(0,1/2]$, let $G$ be a simple graph on $n$ vertices with at least $\delta n^2$ edges.
Is there an absolute constant $c>0$ such that one can always find subgraphs $H_1,H_2\subseteq G$ satisfying:

(1) $e(H_1)\ge c\,\delta^3 n^2$, and for every two edges $e,f\in E(H_1)$ there exists a (simple) cycle in $H_1$ of length at most $6$ containing both $e$ and $f$; moreover if $e$ and $f$ share a vertex then there is a $4$-cycle in $H_1$ containing both.

(2) $e(H_2)\ge c\,\delta^2 n^2$, and for every two edges $e,f\in E(H_2)$ there exists a (simple) cycle in $H_2$ of length at most $8$ containing both $e$ and $f$.

Here $e(H)$ denotes the number of edges.
